% ============================================================
% Lecture deck — Module 8 / Notebook 29
% Architecture Patterns for AI Applications
% HSBI house style (Prof. Dr. Christoph Weisser). Thin deck —
% the notebook is the source of truth.
% ============================================================
\documentclass[aspectratio=169,11pt]{beamer}

\usepackage[utf8]{inputenc}
\usepackage[T1]{fontenc}
\usepackage{lmodern}
\usepackage[english]{babel}
\usepackage{graphicx}
\usepackage{booktabs}
\usepackage{amsmath,amssymb}
\usepackage{xcolor}

\usetheme{Madrid}
\usecolortheme{default}
\setbeamertemplate{navigation symbols}{}
\setbeamertemplate{footline}[frame number]
\setbeamertemplate{blocks}[rounded][shadow=false]
\setbeamertemplate{headline}{%
  \leavevmode%
  \hbox{%
    \begin{beamercolorbox}[wd=\paperwidth,ht=2.5ex,dp=1.125ex]{section in head/foot}%
      \insertsectionnavigationhorizontal{\paperwidth}{}{\hskip0pt plus1filll}%
    \end{beamercolorbox}%
  }%
}
\AtBeginSection[]{%
  \begin{frame}{Outline}
    \tableofcontents[currentsection,sectionstyle=show/shaded,subsectionstyle=show/show/hide]
  \end{frame}
}

\title{Architecture Patterns for AI Applications}
\subtitle{Module 8 --- Business AI in Practice $\cdot$ Notebook 29}
\author{Prof.\ Dr.\ Christoph Weisser}
\institute{HSBI}
\date{Summer Semester 2026}

\begin{document}

\begin{frame}\titlepage\end{frame}

\begin{frame}{Contents}
  \tableofcontents[hideallsubsections]
\end{frame}

% ============================================================
\section{Traffic}
% ============================================================
\begin{frame}{Architecture is the answer to a \emph{traffic} question}
  \begin{block}{Ask three things before you choose}
    \begin{itemize}
      \item \textbf{How many users?} \quad 1 $\to$ 10 $\to$ 100 $\to$ 1{,}000+.
      \item \textbf{How often, and who invokes it?} A human at a keyboard, a schedule, or another service?
      \item \textbf{How many teams} will own and ship pieces of it?
    \end{itemize}
  \end{block}
  \begin{exampleblock}{The point}
    Architecture is not about sophistication --- it is about matching structure to load. The same feature wants a different shape at 5 users than at 5{,}000.
  \end{exampleblock}
\end{frame}

% ============================================================
\section{Patterns}
% ============================================================
\begin{frame}{The four software patterns}
  \small
  \begin{center}
  \begin{tabular}{lp{4.0cm}p{3.6cm}}
    \toprule
    \textbf{Pattern} & \textbf{When it fits} & \textbf{Example} \\
    \midrule
    Single-tier        & one analyst exploring, or one daily job & Jupyter notebook; cron + script \\
    Three-tier         & interactive use by non-developers & Streamlit $+$ FastAPI $+$ DB \\
    Service-oriented   & 2--3 capabilities, one team & shared infra, common interface \\
    Microservices      & 4+ teams shipping independently & one team per service \\
    \bottomrule
  \end{tabular}
  \end{center}
  \vspace{4pt}
  Read top to bottom as growing traffic and growing team count --- not as ``primitive to advanced''.
\end{frame}

\begin{frame}{The over-engineering trap}
  \begin{alertblock}{Microservices for a 5-user pilot}
    Microservices became fashionable around 2015; many teams adopted them \emph{before} they had the team size or traffic to justify the operational cost. The right unit of a service is a \emph{capability} --- not every Python module.
  \end{alertblock}
  \begin{exampleblock}{The rule}
    Pick the architecture you'll have \emph{outgrown in 12 months}, not the one you'll \emph{need in 36}. Migrations are real, but rarely the thing that kills a project --- premature complexity is.
  \end{exampleblock}
\end{frame}

% ============================================================
\section{Pipeline}
% ============================================================
\begin{frame}{The cross-cutting ML / AI pipeline}
  \begin{block}{What pure-software patterns don't capture}
    \begin{itemize}
      \item \textbf{Training data is part of the artefact.} Two models that look identical but trained on different snapshots are \emph{different products}.
      \item \textbf{Drift exists.} 92\,\% accurate in March is not 92\,\% in October --- monitoring must compare against fresh ground truth.
      \item \textbf{The feedback loop has economic cost.} Every prediction nudges the world the next training batch will learn from.
    \end{itemize}
  \end{block}
  Data $\to$ features $\to$ train $\to$ evaluate $\to$ serve $\to$ \emph{monitor} $\to$ (back to data).
\end{frame}

% ============================================================
\section{Choosing}
% ============================================================
\begin{frame}{Choosing the right size}
  \begin{block}{Smells that you've outgrown your current shape}
    \begin{itemize}
      \item Someone asks ``can we run this nightly?'' $\to$ add scheduling (single-tier job).
      \item A non-developer needs to use it $\to$ add a UI tier (three-tier).
      \item A second team wants to build on it $\to$ expose it as a service.
    \end{itemize}
  \end{block}
  \begin{exampleblock}{Two opposite mistakes to avoid}
    \emph{Over-engineering} (microservices for a pilot) and \emph{under-engineering} (a 200-user system held together by one notebook).
  \end{exampleblock}
\end{frame}

% ============================================================
\section{Course map}
% ============================================================
\begin{frame}{How the course maps to these patterns}
  \small
  \begin{center}
  \begin{tabular}{ll}
    \toprule
    \textbf{Pattern} & \textbf{Where it appears in the course} \\
    \midrule
    Single-tier (script / notebook) & Every notebook NB 1--27, by default \\
    Scheduled single-tier           & NB 24 --- scheduling \& orchestration \\
    Three-tier with AI in the middle& NB 34 (POC 2) and NB 27 capstone \\
    Service-oriented                & \texttt{llm\_providers.py} --- pluggable providers \\
    Microservices                   & \emph{patterns only} --- too big for a teaching course \\
    ML pipeline                     & NB 15--17 $+$ NB 24 --- train / evaluate / serve \\
    \bottomrule
  \end{tabular}
  \end{center}
  \vspace{4pt}
  You \emph{build} the three-tier and ML-pipeline shapes in \textbf{Notebook 34}; the rest you meet as patterns here.
\end{frame}

\end{document}
